\section{Plan de trabajo}
\label{sec:planTrabajo}

El desarrollo del proyecto ha seguido un enfoque incremental e iterativo, condicionado por la naturaleza encadenada del pipeline. Esto se pone de manifiesto en la \hyperref[fig:arquitectura-software]{Figura~\ref*{fig:arquitectura-software}}, donde se representa un esquema del pipeline y se indican las distintas fases del trabajo y sus módulos.  Cada módulo utiliza como entrada la salida de etapas anteriores, por lo que los errores producidos en una fase pueden propagarse y afectar al rendimiento de las siguientes. Por este motivo, el avance del sistema no fue estrictamente lineal, sino que requirió revisar y ajustar de forma continua módulos ya desarrollados a medida que se incorporaban nuevas funcionalidades.

La primera fase del trabajo se centró en la detección de los elementos principales del juego. Para ello, se partió de un modelo YOLO entrenado con datos etiquetados para detectar jugadores, porteros, árbitros y balón. A partir de esta base, se incorporaron progresivamente la calibración del campo mediante homografía, apoyada en PnLCalib, la identificación automática de equipos a partir del color de la camiseta, el tracking multiobjeto mediante una adaptación de ByteTrack y una capa de tracking canónico destinada a mantener identidades más estables a lo largo del vídeo. Estas tareas se desarrollaron de forma conjunta y fueron evaluadas tanto con métricas como visualizando los resultados sobre vídeos anotados y permitiendo asi la corrección de problemas como pérdidas de identidad, intercambios entre jugadores, detecciones inestables del balón o errores en la asignación de equipo.

Una vez obtenidas trayectorias estables, se abordó una segunda fase orientada al análisis táctico del partido. En esta etapa se construyó una representación estructurada de la información generada por el bloque de tracking, incluyendo identidades, equipos, posiciones en el campo y evolución temporal de los jugadores y del balón. Sobre esta representación se incorporó la estimación de la posesión y la clasificación posicional de los jugadores mediante un modelo Set Transformer entrenado con datos etiquetados manualmente. Esta fase también requirió varias iteraciones, ya que la calidad del análisis dependía directamente de la estabilidad de las detecciones, del tracking y de la información geométrica disponible.

La tercera fase se centró en la narración automática del partido y se planteó como un bloque asíncrono respecto al procesamiento frame a frame, como se refleja en la \hyperref[fig:arquitectura-software]{Figura~\ref*{fig:arquitectura-software}}. A partir de la información estructurada del partido, se adaptó la salida del tracking al formato requerido por PathCRF para detectar acciones relevantes, como pases, controles o tiros. Estos eventos se utilizaron posteriormente como entrada para un módulo de generación de comentarios mediante un LLM local, junto con información contextual del partido. Para ello fue necesario definir qué información debía recibir el modelo, cómo controlar la longitud y variabilidad de los comentarios y cómo incorporar comentarios contextuales basados en estadísticas, equipos o datos de la competición.

Finalmente, los comentarios generados se transformaron en audio mediante un módulo de síntesis de voz en el que se evaluaron varias alternativas, incluyendo XTTS, Qwen TTS y ElevenLabs, contemplando además la alternancia de voces disponibles para producir una narración más dinámica.

Una vez integrado el pipeline completo, se evaluó el sistema atendiendo a métricas de detección, tracking, cobertura de jugadores y balón, clasificación posicional, generación de eventos y viabilidad de ejecución local. Además, se revisaron los distintos módulos del sistema con el objetivo de identificar oportunidades de optimización y analizar el compromiso entre precisión, calidad de los resultados y latencia.

Esta evaluación permitió identificar las principales limitaciones del sistema, especialmente en relación con la detección del balón, las oclusiones, la estabilidad de identidades y el coste computacional de ejecutar el pipeline completo en un entorno local con una GPU RTX 3080.

Todo el software del proyecto que se ha desarrollado puede encontrarse en el siguiente repositorio público: \url{https://github.com/hectorgarciaa/narrador-futbol.git }


\begin{figure}[H]
\centering
\resizebox{0.95\textwidth}{!}{%
\begin{tikzpicture}[
    font=\scriptsize,
    node distance=0.45cm and 0.7cm,
    block/.style={
        rectangle,
        draw=black,
        rounded corners=2pt,
        minimum width=4.0cm,
        minimum height=0.72cm,
        align=center,
        fill=white
    },
    title/.style={
        font=\scriptsize\bfseries,
        fill=white,
        inner sep=1.5pt
    },
    container/.style={
        rectangle,
        draw=black,
        rounded corners=4pt,
        inner sep=8pt,
        fill=none
    },
    subcontainer/.style={
        rectangle,
        draw=black!70,
        dashed,
        rounded corners=4pt,
        inner sep=6pt,
        fill=none
    },
    arrow/.style={-{Latex}, thick}
]

% =========================================================
% BLOQUE SÍNCRONO: FASE 1 Y FASE 2
% =========================================================

\node[title] (fase1label) {Fase 1 --- Tracking y localización};

\node[block, below=0.25cm of fase1label] (det)
{Detección (YOLO)};

\node[block, below=of det] (hom)
{Homografía (PnLCalib)};

\node[block, below=of hom] (equipos)
{Identificación de equipos};

\node[block, below=of equipos] (byte)
{ByteTrack};

\node[block, below=of byte] (canonico)
{Tracking canónico};

\node[title, below=0.55cm of canonico] (fase2label)
{Fase 2 --- Análisis táctico};

\node[block, below=0.25cm of fase2label] (posesion)
{Posesión};

\node[block, below=of posesion] (posiciones)
{Clasificación posicional\\(Set Transformer)};

\node[inner sep=0pt, minimum size=0pt] (syncpad)
at ($(fase1label.north)+(0,0.45)$) {};

\node[subcontainer,
      fit=(syncpad)(fase1label)(det)(hom)(equipos)(byte)(canonico)
          (fase2label)(posesion)(posiciones)] (syncbox) {};

\node[title, anchor=north west]
at ($(syncbox.north west)+(0.15,-0.05)$)
{Bloque síncrono};

% =========================================================
% BLOQUE ASÍNCRONO: FASE 3
% =========================================================

\coordinate (asynccenter) at ($(syncbox.east)+(4.8cm,0)$);

\node[block] (comentario) at (asynccenter)
{Generación de comentario\\(LLM local)};

\node[block, above=of comentario] (acciones)
{Detección de acciones\\(PathCRF)};

\node[block, below=of comentario] (tts)
{Síntesis de voz\\(XTTS)};

\node[title, above=0.25cm of acciones] (fase3label)
{Fase 3 --- Narración automática};

\node[inner sep=0pt, minimum size=0pt] (asyncpad)
at ($(fase3label.north)+(0,0.45)$) {};

\node[subcontainer,
      fit=(asyncpad)(fase3label)(acciones)(comentario)(tts)] (asyncbox) {};

\node[title, anchor=north west]
at ($(asyncbox.north west)+(0.15,-0.05)$)
{Bloque asíncrono};

% =========================================================
% CONTENEDOR ÚNICO: PIPELINE
% =========================================================

\node[inner sep=0pt, minimum size=0pt] (pipepad)
at ($(syncbox.north west)!0.5!(asyncbox.north east)+(0,0.95)$) {};

\node[container, fit=(pipepad)(syncbox)(asyncbox)] (pipelinebox) {};

\node[title, anchor=north west]
at ($(pipelinebox.north west)+(0.18,-0.08)$)
{Pipeline};

% =========================================================
% NODO EXTERNO: VÍDEO BROADCAST
% =========================================================

\node[block, anchor=south] (video)
at ($(pipelinebox.north)+(0,0.65cm)$)
{Vídeo broadcast};

% =========================================================
% FLECHAS
% =========================================================

% Entrada externa hacia el pipeline:
% baja y gira una sola vez hacia la izquierda hasta entrar por det.east
\draw[arrow] (video.south) |- (det.east);

% Flujo síncrono
\draw[arrow] (det) -- (hom);
\draw[arrow] (hom) -- (equipos);
\draw[arrow] (equipos) -- (byte);
\draw[arrow] (byte) -- (canonico);

\draw[arrow] (canonico) -- (fase2label);
\draw[arrow] (fase2label) -- (posesion);
\draw[arrow] (posesion) -- (posiciones);

% Paso del bloque síncrono al bloque asíncrono
\draw[arrow] (posiciones.east) -- ++(0.7,0) |- (acciones.west);

% Flujo asíncrono
\draw[arrow] (acciones) -- (comentario);
\draw[arrow] (comentario) -- (tts);

\end{tikzpicture}
}
\caption[Arquitectura software del pipeline]{Arquitectura software del pipeline del sistema. Se organiza en un bloque síncrono, que agrupa las fases de detección, localización, tracking y análisis táctico, y un bloque asíncrono, responsable de la narración automática.}
\label{fig:arquitectura-software}
\end{figure}
